\subsection{Fresh composition and cross-potential confirmation}\label{sec:fresh-physics}
After the earlier source and physics studies, we freeze the Gaussian model,
LatentSpring's harmonic model, and both fixed-coefficient physics-update variants.
The cheaper unweighted-force update is designated the primary physical candidate
before querying new outcomes. We sample a new pool of 256 distinct compositions
from the original development partition, excluding every old official candidate
composition. Reference qualification accepts 121; fixed metadata ranks select
8 per 17--20, 21--24 and 25--28 atom bin, yielding 24 new compositions. All selected
compositions are absent from both verified processed corpora. The reference
coordinates do not enter generation or fitting. Every model generates 32 samples
per composition in each of the same two existing continuations: 6,144 new outputs,
with no additional model fitting or parameter selection. Inference uses the
network alone: no per-output energy ranking or coordinate optimization is applied.

\begin{table}[t]
\centering
\caption{Fresh 24-composition graph support, 768 attempts per continuation.
All rates describe raw outputs before any structural relaxation. All methods use
128 primitive denoiser calls per attempted output.}
\begin{tabular}{lrrr}
\toprule
Model & Continuation 1 & Continuation 2 & Pooled rate\\
\midrule
Gaussian reference &332&352&44.53\%\\
LatentSpring harmonic &435&375&52.73\%\\
With force update &439&382&53.45\%\\
With complete-work update &426&383&52.67\%\\
\bottomrule
\end{tabular}\label{tab:fresh-graph}
\end{table}

The harmonic source improves graph support over Gaussian by 8.20 percentage points,
with conditional paired 95\% interval $[5.08,11.33]$. Both continuation signs are
positive. The primary force update has an observed graph-rate difference of
$+0.72$ points versus harmonic, with interval $[-0.98,2.34]$. This is compatible
with unchanged validity and excludes a decline exceeding two points under this
conditional interval; it is not evidence of a statistically significant validity
increase or a population-wide noninferiority guarantee. All accepted graphs have
distinct perceived connectivity within each composition in these draws.

\paragraph{Independent physical readout.}
We evaluate every generated coordinate set with eSEN and with gas-phase GFN2-xTB
single points and analytic gradients \citep{gfn2xtb}. GFN2 is not used for training
or method selection. It provides a different approximate physical model, not
DFT ground truth. We retain all failed calculations: 4 of 6,144 generated-geometry
GFN2 calculations fail, all outside graph support. Every graph-supported output
succeeds, and all original-reference inversion checks pass. Total physical cost
is 12,384 raw eSEN queries and 6,240 GFN2 attempts, including reference checks.
The audited eSEN readout uses $E_+(X)=[E(X)+E(-X)]/2$ and
$F_+(X)=[F(X)-F(-X)]/2$; force RMS is the root mean squared atom-force norm.
GFN2 inputs are fixed at their saved generated coordinates.
The xTB command uses \texttt{--grad}, never \texttt{--opt}; gradients are readouts
and do not update coordinates. Physical teacher displacements occur only in the
earlier training procedure, on separate FIT samples.

\begin{table}[t]
\centering
\caption{Primary force-update change versus the frozen harmonic model on paired
common-graph support. Both continuation signs improve for both readouts under
both potentials. Intervals are conditional paired 95\% intervals.}
\small
\setlength{\tabcolsep}{4pt}
\begin{tabular}{lrr}
\toprule
Potential & Energy change (eV/atom) & Force-RMS change (eV/\AA)\\
\midrule
eSEN & $-0.01308\;[-0.01805,-0.00817]$ & $-0.34793\;[-0.41498,-0.28478]$\\
GFN2 & $-0.01168\;[-0.01616,-0.00724]$ & $-0.35620\;[-0.42983,-0.28714]$\\
\bottomrule
\end{tabular}\label{tab:fresh-physical}
\end{table}

All 24 composition cells in each continuation have comparable graph-supported
pairs for the primary contrast. The prespecified cross-potential confirmation
gate passes, including success coverage and the observed two-point validity
criterion. All-attempt eSEN and common-success GFN2 comparisons agree in direction;
selected-support results remain distinct from unconditional or equilibrium
expectations. Source-only harmonic-versus-Gaussian energy improvements are pooled
positive but heterogeneous across continuations, so we do not describe them as
repeated physical-quality gains in both models.

Complete-work versus force-update energy differences remain unresolved:
$-0.00413$ eV/atom under eSEN, interval $[-0.00903,0.00078]$, and $-0.00363$ under
GFN2, interval $[-0.00792,0.00075]$, on common graph support. Force and graph
intervals also span zero. Thus the fresh result supports the practical physical
feedback update, while not establishing an extra advantage of Jarzynski weighting.
The previously defined local work identity remains mathematically useful, but
these results do not imply globally Boltzmann-distributed generator outputs.

